% Part: normal-modal-logic
% Chapter: filtrations
% Section: more-filtrations

\documentclass[../../../include/open-logic-section]{subfiles}

\begin{document}

\olfileid{nml}{fil}{acc}

\olsection{الترشيحات وخواص النفاذ}

كما أشرنا، تكون ترشيحات النماذج الكلية والمتسلسلة والانعكاسية كلية أو
متسلسلة أو انعكاسية دائمًا. لكن ليس كل ترشيح لنموذج متناظر أو متعدٍ متناظرًا
أو متعديًا، على الترتيب. ومع ذلك، يمكن في بعض الحالات تعريف الترشيحات بحيث
يصح ذلك. ولتحقيقه، نسير كما في تعريف الترشيح الأخشن، لكننا نضيف شروطًا أخرى
إلى تعريف~$R^*$. لتكن~$\Gamma$ مغلقة تحت أخذ ال!!{formula}s الجزئية.
انظر إلى العلاقات~$C_i(u,v)$ في~\olref{tab:Cn-filtrations} بين
العالمين~$u$ و$v$ في نموذج~$\mModel{M}=\tuple{W,R,V}$. ويمكننا تعريف
$R^*[u][v]$ على أساس توليفات من هذه الشروط. فمثلًا، إذا اشترطنا أن
$R^*[u][v]$ إذا وفقط إذا تحقق الشرط~$C_1(u,v)$، حصلنا بالضبط على
الترشيح الأخشن. وإذا اشترطنا أن~$R^*[u][v]$ إذا وفقط إذا تحقق الشرطان
$C_1(u,v)$ و$C_2(u,v)$ معًا، حصلنا على ترشيح مختلف. وهو «أدق» من
الترشيح الأخشن، لأن عدد أزواج العوالم التي تستوفي~$C_1(u,v)$ و$C_2(u,v)$
أقل من عدد الأزواج التي تستوفي~$C_1(u,v)$ وحده.

\begin{table}[ht]
  \centering
  \begin{tabular}{|ll|}
    \hline
    \multirow{2}{*}{$C_1(u,v)$:}
    \iftag{prvBox}{&
      إذا كان~$\Box!A \in \Gamma$ و$\mSat{M}{\Box!A}[u]$،
      كان~$\mSat{M}{!A}[v]$؛\iftag{prvDiamond}{ و}{}\\}{}
    \iftag{prvDiamond}{&
      إذا كان~$\Diamond!A \in \Gamma$ و$\mSat{M}{!A}[v]$،
      كان~$\mSat{M}{\Diamond!A}[u]$؛ \\}{}
    \hline
    \multirow{2}{*}{$C_2(u,v)$:}
    \iftag{prvBox}{&
      إذا كان~$\Box!A \in \Gamma$ و$\mSat{M}{\Box!A}[v]$،
      كان~$\mSat{M}{!A}[u]$؛\iftag{prvDiamond}{ و}{}\\}{}
    \iftag{prvDiamond}{&
      إذا كان~$\Diamond!A \in \Gamma$ و$\mSat{M}{!A}[u]$،
      كان~$\mSat{M}{\Diamond!A}[v]$؛ \\}{}
    \hline
    \multirow{2}{*}{$C_3(u,v)$:}
    \iftag{prvBox}{&
      إذا كان~$\Box!A \in \Gamma$ و$\mSat{M}{\Box!A}[u]$،
      كان~$\mSat{M}{\Box!A}[v]$؛\iftag{prvDiamond}{ و}{}\\}{}
    \iftag{prvDiamond}{&
      إذا كان~$\Diamond!A \in \Gamma$ و$\mSat{M}{\Diamond!A}[v]$،
      كان~$\mSat{M}{\Diamond!A}[u]$؛ \\}{}
    \hline
    \multirow{2}{*}{$C_4(u,v)$:}
    \iftag{prvBox}{&
      إذا كان~$\Box!A \in \Gamma$ و$\mSat{M}{\Box!A}[v]$،
      كان~$\mSat{M}{\Box!A}[u]$؛\iftag{prvDiamond}{ و}{}\\}{}
    \iftag{prvDiamond}{&
      إذا كان~$\Diamond!A \in \Gamma$ و$\mSat{M}{\Diamond!A}[u]$،
      كان~$\mSat{M}{\Diamond!A}[v]$؛ \\}{}
    \hline
    \end{tabular}
  \caption{شروط على العوالم الممكنة لتعريف الترشيحات.}
  \ollabel{tab:Cn-filtrations}
\end{table}

\begin{thm}\ollabel{thm:more-filtrations}
  ليكن~$\mModel{M}=\tuple{W,R,V}$ نموذجًا، ولتكن~$\Gamma$ مغلقة تحت
  أخذ ال!!{formula}s الجزئية. ولتُعرّف~$W^*$ و$V^*$ كما في
  \olref[fil]{defn:filtration}. عندئذ:
  \begin{enumerate}
  \item افترض أن~$R^*[u][v]$ إذا وفقط إذا كان
    $C_1(u,v)\land C_2(u,v)$. عندئذ تكون~$R^*$ متناظرة، ويكون
    $\mModel{M^*}=\tuple{W^*,R^*,V^*}$ ترشيحًا إذا كان~$\mModel{M}$
    متناظرًا.
  \item افترض أن~$R^*[u][v]$ إذا وفقط إذا كان
    $C_1(u,v)\land C_3(u,v)$. عندئذ تكون~$R^*$ متعدية، ويكون
    $\mModel{M^*}=\tuple{W^*,R^*,V^*}$ ترشيحًا إذا كان~$\mModel{M}$
    متعديًا.
  \item افترض أن~$R^*[u][v]$ إذا وفقط إذا كان
    $C_1(u,v)\land C_2(u,v)\land C_3(u,v)\land C_4(u,v)$. عندئذ تكون
    $R^*$ متناظرة ومتعدية، ويكون~$\mModel{M^*}=\tuple{W^*,R^*,V^*}$
    ترشيحًا إذا كان~$\mModel{M}$ متناظرًا ومتعديًا.
  \item افترض أن~$R^*$ معرّفة بحيث يكون~$R^*[u][v]$ إذا وفقط إذا كان
    $C_1(u,v)\land C_3(u,v)\land C_4(u,v)$. عندئذ تكون~$R^*$ متعدية
    وإقليدية، ويكون~$\mModel{M^*}=\tuple{W^*,R^*,V^*}$ ترشيحًا إذا
    كان~$\mModel{M}$ متعديًا وإقليديًا.
  \end{enumerate}
\end{thm}

\begin{proof}
  \begin{enumerate}
    \item يتضح فورًا أن~$R^*$ متناظرة، لأن
      $C_1(u,v)\Leftrightarrow C_2(v,u)$ و
      $C_2(u,v)\Leftrightarrow C_1(v,u)$. فيبقى أن نبين أنه إذا كان
      $\mModel{M}$ متناظرًا، كان~$\mModel{M^*}$ ترشيحًا عبر~$\Gamma$.
      ويضمن الشرط~$C_1(u,v)$ استيفاء
      \iftag{prvBox}{%
        \iftag{prvDiamond}{%
          \olref[fil]{defn:filtration-R2} و
          \olref[fil]{defn:filtration-R3} من
          \olref[fil]{defn:filtration}}{%
          \olref[fil]{defn:filtration}\olref[fil]{defn:filtration-R2}}}{%
        \olref[fil]{defn:filtration}\olref[fil]{defn:filtration-R3}}.
      فلا يبقى إلا التحقق من
      \olref[fil]{defn:filtration}\olref[fil]{defn:filtration-R1}، أي
      إن~$Ruv$ يستلزم~$R^*[u][v]$.

      افترض إذن~$Ruv$. ولإثبات~$R^*[u][v]$، علينا إثبات~$C_1(u,v)$
      و$C_2(u,v)$. أما~$C_1$: \iftag{prvBox}{فإذا كان~$\Box!A
        \in\Gamma$ و$\mSat{M}{\Box!A}[u]$، كان أيضًا
        $\mSat{M}{!A}[v]$ (لأن~$Ruv$).\iftag{prvDiamond}{ وبالمثل،
        }{}}{}\iftag{prvDiamond}{إذا كان~$\Diamond!A \in \Gamma$
        و$\mSat{M}{!A}[v]$، كان~$\mSat{M}{\Diamond!A}[u]$ لأن
        $Ruv$.}{} وأما~$C_2$: \iftag{prvBox}{فإذا كان~$\Box!A \in
        \Gamma$ و$\mSat{M}{\Box!A}[v]$، استلزم~$Ruv$ أن~$Rvu$
        بالتناظر، ومن ثم~$\mSat{M}{!A}[u]$.\iftag{prvDiamond}{
        وبالمثل، }{}}{}\iftag{prvDiamond}{إذا كان~$\Diamond!A
        \in\Gamma$ و$\mSat{M}{!A}[u]$، كان
        $\mSat{M}{\Diamond!A}[v]$ (لأن~$Rvu$ بالتناظر).}{}
    \item تمرين.
    \item تمرين.
    \item تمرين.
  \end{enumerate}
\end{proof}

\begin{prob}
  أكمل برهان~\olref[nml][fil][acc]{thm:more-filtrations}.
\end{prob}

\end{document}
